\begin{exercise}
\label{ex:second-incompleteness-01}
Reprove \cref{thm:godel-second} directly from a G\"odel sentence by showing in
\(T\) that \(\mathsf{Con}(T)\to G_T\).  Identify the derivability conditions
used in the internal argument.
\end{exercise}

\begin{exercise}
\label{ex:second-incompleteness-02}
Let \(T\) satisfy the hypotheses of L\"ob's theorem.  Prove that if
\(T\derives\mathsf{Prov}_T(\ulcorner\varphi\urcorner)
\leftrightarrow\varphi\), then \(T\derives\varphi\).
\end{exercise}

\begin{exercise}
\label{ex:second-incompleteness-03}
Formulate a syntactic version of Tarski's theorem: show that no consistent
theory extending \(\mathsf Q\) can prove every biconditional
\[
  \mathsf{Tr}(\overline{\ulcorner\sigma\urcorner})
  \leftrightarrow\sigma
\]
for all arithmetic sentences \(\sigma\).  Use the diagonal lemma and state
where consistency is required.
\end{exercise}

\begin{exercise}
\label{ex:second-incompleteness-04}
Audit the sentence ``G\"odel proved that true statements cannot be proved.''
Replace it by three distinct correct statements concerning \(G_T\),
\(\mathsf{Con}(T)\), and \(\Th(\mathcal N)\), each with exact hypotheses.
\end{exercise}
